% =====================================================================
%  Wspólna preambuła szablonów pism ZUS (świadczenia rodzicielskie 2026)
%  Silnik: XeLaTeX albo Tectonic (pdflatex NIE obsłuży polskich znaków w fontspec).
%  Kompilacja:  tectonic -X compile plik.tex
% =====================================================================
\documentclass[11pt,a4paper]{article}

\usepackage{fontspec}
\usepackage{polyglossia}
\setmainlanguage{polish}
\setotherlanguage{english}
% font domyślny (Latin Modern) — Tectonic go dostarcza

\usepackage[margin=2.5cm,a4paper]{geometry}
\usepackage{graphicx}
\usepackage{xcolor}
\usepackage{array}
\usepackage{enumitem}
\usepackage{fancyhdr}
\usepackage{hyperref}

\definecolor{zus}{HTML}{007A33}      % zieleń ZUS
\definecolor{szary}{HTML}{6B7280}
\definecolor{ramka}{HTML}{9CA3AF}

\hypersetup{colorlinks=true, linkcolor=zus, urlcolor=zus, pdftitle={Szablony pism ZUS}}

% --- pole do wypełnienia odręcznego -----------------------------------
% \Pole{nazwa} rysuje kropkowaną linię z podpisem pola pod spodem.
% Pole do wypełnienia: kropkowana linia, pod nią drobny podpis pola.
\newcommand{\Pole}[1]{%
  \textcolor{szary}{\dotfill}\,%
  {\scriptsize\textcolor{szary}{\textsc{#1}}}%
}

% --- blok warunkowy: treść pokazywana tylko, gdy pole prawdziwe --------
\newenvironment{warunek}{%
  \par\smallskip\noindent\begin{minipage}{\linewidth}%
  \color{black}%
}{%
  \end{minipage}\par\smallskip%
}

% --- nagłówek pisma ---------------------------------------------------
\newcommand{\NaglowekPisma}[1]{%
  \par\vspace{0.8em}%
  \begin{center}%
    {\large\bfseries\MakeUppercase{#1}}%
  \end{center}%
  \vspace{0.4em}%
}

% --- blok nadawcy (pierwsze linie pisma) -------------------------------
\newenvironment{nadawca}{%
  \begin{flushleft}\small
}{%
  \end{flushleft}\normalsize
}

% --- stopka z numerem strony -----------------------------------------
\pagestyle{fancy}
\fancyhf{}
\renewcommand{\headrulewidth}{0pt}
\fancyfoot[C]{\textcolor{szary}{\small\thepage}}

% --- metadane szablonu na początku pisma ------------------------------
\newcommand{\Metadane}[3]{%
  \begin{flushright}\footnotesize\textcolor{szary}%
    \textbf{#1}\\[2pt] {\itshape #2}\\[2pt] #3%
  \end{flushright}\vspace{0.5em}%
}

\setlength{\parindent}{0pt}
\setlength{\parskip}{0.6em}

\begin{document}

\begin{flushright}
{\footnotesize\color{szary}
\textbf{Szablon P16}\\
{\itshape Wniosek o uzupełniający urlop macierzyński (wcześniak / szpital)}\\
art. 180\textsuperscript{1} k.p.; art. 29 ust. 5 ustawy zasiłkowej -- Formularz: ZAM}\
\end{flushright}

\vspace{0.5em}
{\color{ramka}\hrule height 0.4pt}
\vspace{0.8em}

\NaglowekPisma{WNIOSEK O UDZIELENIE UZUPEŁNIAJĄCEGO URLOPU MACIERZYŃSKIEGO}

Na podstawie art. 183\textsuperscript{1} \S{} 1 pkt 1 Kodeksu pracy wnoszę o udzielenie mi
uzupełniającego urlopu macierzyńskiego w wymiarze stanowiącym różnicę
między 41 tygodniami a wymiarami urlopów wykorzystanych, w okresie
od \Pole{urlop od} do \Pole{urlop do}.

Uzupełniający urlop macierzyński przysługuje w wymiarze:

dzień porodu do 28. tygodnia ciąży        = 41 tygodni − wykorzystane tygodnie
od 28. do 36. tygodnia ciąży             = 41 tygodni − wykorzystane tygodnie

Tydzień ciąży w chwili porodu: \Pole{tydzien ciazy}
Masa urodzeniowa dziecka: \Pole{masa} g
Pobyt dziecka w szpitalu: od \Pole{szpital od} do \Pole{szpital do}  (\Pole{dni hospitalizacji} dni)

\Pole{data pisma}
.....................................
(podpis)

\vfill
{\footnotesize\textcolor{szary}\itshape Uzupełniający urlop macierzyński przysługuje w razie urodzenia wcześniaka albo gdy dziecko wymagało hospitalizacji po porodzie (art. 183\textsuperscript{1} \S{} 1 pkt 1--2 k.p.). }

\end{document}
